\section{Finite quotients of prescribed excess order}\label{finite-quotients-of-prescribed-excess-order}

The next theorem separates a universal finite search space from the task
of proving that a particular family has the required class-two quotient.
That identification is an application hypothesis, not a consequence of
the search.

\textbf{Theorem UA.Nilpotent (complete marked kernel parameter space).}
Let \(F=\langle x_i:i\in I\rangle\), where \(|I|=d\) and \(x_i^2=1\).
Suppose \(F\) has all Boolean deletions. Put \(Q=F/\gamma_3(F)\), with
its image marking. For \(k\geq0\), set \[
U=F/\gamma_{k+3}(F),\qquad K=\gamma_3(U),\qquad
K_S=K\cap U_S.
\tag{N1}
\] Then \(U\) is a finite two-group, \[
|U|\leq 2^{\,\sum_{j=1}^{k+2}d^j},
\tag{N2}
\] the markings of \(U\) and \(Q\) have Boolean deletions and IP, and
\(U/K\) is canonically the marked group \(Q\). For every
\(S\subseteq I\), \[
K_S=\gamma_3(U_S).
\tag{N3}
\]

There is a bijection between the following two sets:

\begin{enumerate}
\def\labelenumi{\arabic{enumi}.}
\tightlist
\item
  Ordered marked isomorphism classes of finite two-group quotients
  \(f:F\twoheadrightarrow G\) for which the induced marked map
  \(Q\to G/\gamma_3(G)\) is an isomorphism, \(|G|=|Q|2^k\), and \(G\)
  has IP.
\item
  Normal subgroups \(N\unlhd U\) satisfying \(N\leq K\), \([K:N]=2^k\),
  and \[
  \frac{NK_S}{N}\cap\frac{NK_T}{N}
  =\frac{NK_{S\cap T}}{N}
  \quad(S,T\subseteq I).
  \tag{N4}
  \]
\end{enumerate}

The correspondence sends \(N\) to \(U/N\) with the original image
markers. If \(D\leq\operatorname{Sym}(I)\) is specified and each
\(\sigma\in D\) induces an automorphism \(x_i\mapsto x_{\sigma(i)}\) of
\(F\), then equivalence allowing precisely these label permutations is
the \(D\)-orbit relation on the kernels in item 2.

\emph{Proof.} The degree-one part of the lower-central graded Lie ring
of \(F\) is generated by elements of exponent dividing two. By
multilinearity, every iterated commutator in those generators also has
exponent dividing two in its homogeneous component. In degree \(j\), at
most \(d^j\) left-normed commutators suffice to generate that component.
Consequently \[
|\gamma_j(F)/\gamma_{j+1}(F)|\leq2^{d^j}.
\] Multiplying these bounds for \(1\leq j\leq k+2\) proves finiteness,
the two-power order and (N2). Each deletion preserves every
lower-central term, and therefore descends both to \(U\) and to \(Q\).
Proposition UA.D gives IP for both image markings; no separate IP
hypothesis on \(Q\) is needed. The epimorphic-image property of the
lower central series gives \(U/K\cong F/\gamma_3(F)=Q\).

For (N3), one inclusion follows from \(U_S\leq U\). If
\(z\in\gamma_3(U)\cap U_S\), the deletion \(e_S\) fixes \(z\) and sends
\(\gamma_3(U)\) into \(\gamma_3(U_S)\). This proves the reverse
inclusion.

Let \(G\) satisfy item 1. Its subgroup \(\gamma_3(G)\) has order
\(2^k\). In a nilpotent group, equality between two successive
nontrivial lower-central terms would persist at every later term,
contradicting termination of that series. Thus each nontrivial step
beginning at \(\gamma_3(G)\) strictly decreases its order. There are at
most \(k\) such steps, and \(\gamma_{k+3}(G)=1\). The map \(F\to G\)
therefore factors through \(U\). Its kernel \(N\) lies in \(K\), since
the induced map \(U/K\to G/\gamma_3(G)\) is an isomorphism, and the
orders give \([K:N]=2^k\).

Conversely, let \(N\unlhd U\) with \(N\leq K\) and \([K:N]=2^k\), and
write \(G=U/N\). Its third lower-central subgroup is \(K/N\), its
canonical class-two quotient is \(Q\), and its order is \(|Q|2^k\). If
\(q:U\to G\), then \[
(K/N)\cap G_S=q(K\cap U_S)=NK_S/N.
\tag{N5}
\] Indeed, \(q(u)\in K/N\) with \(u\in U_S\) implies \(u\in K\), because
\(q^{-1}(K/N)=K\). The reverse inclusion is immediate. Applying UA.K to
\(G\to Q\) proves that IP in \(G\) is equivalent to precisely (N4).

An isomorphism preserving each ordered marker commutes with the two maps
from \(U\), because it agrees with them on a generating family. Their
kernels are consequently equal. Equal kernels plainly give the same
marked quotient. Finally, every specified label permutation acts on
\(F\), preserves its characteristic lower-central terms, and hence acts
on \(U\) and on the displayed kernels. A marker-permuting isomorphism
induces the corresponding kernel equality after this automorphism, and
the converse descends that automorphism to the quotients. This proves
the orbit assertion at exactly the specified equivalence scope.
\(\square\)

\textbf{Theorem UA.Ladder (invariant index-two descent).} Let \(U\) be a
finite two-group and \(M\unlhd U\). The \(U\)-normal subgroups
\(L\leq M\) of index two are exactly the inverse images of the
hyperplanes of \[
M/\bigl(\Phi(M)[M,U]\bigr).
\tag{N6}
\] For \(K\unlhd U\), every \(U\)-normal subgroup \(N\leq K\) of index
\(2^k\) is obtained by \(k\) successive such steps.

\emph{Proof.} An index-two subgroup \(L\) of \(M\) contains \(\Phi(M)\)
and has quotient \(C_2\). If \(L\) is \(U\)-normal, conjugation by \(U\)
on \(M/L\) is trivial, since \(\operatorname{Aut}(C_2)=1\). Thus
\([M,U]\leq L\). Conversely, a hyperplane inverse image has index two
and contains \([M,U]\); for \(l\in L\) and \(u\in U\), the conjugate of
\(l\) differs from \(l\) by an element of \([M,U]\), so lies in \(L\).
The quotient in (N6) is elementary abelian by Burnside's basis theorem,
proving the first assertion.

Suppose \(N\unlhd U\) and \(N<M\). The nontrivial finite two-group
\(M/N\) has a nonzero Frattini quotient \(V\). Conjugation by \(U/N\)
acts on the nonzero linear maps \(V\to\mathbb F_2\). If \(\dim V=a\),
there are \(2^a-1\) such maps, an odd number. Every orbit has two-power
size, so at least one orbit has size one. The kernel of a fixed nonzero
map lifts to a \(U\)-normal subgroup between \(N\) and \(M\), of index
two in \(M\). Repeating with the new subgroup gives a chain terminating
at \(N\); its length is \(k\), since the total index is \(2^k\).
\(\square\)

\textbf{Lemma UA.Order (presentation recognition by a sharp bound).} Let
\(A_1,\ldots,A_t\) be finite normal subgroups of \(G\), with
\(|A_j|\leq a_j\). Suppose \([G:A_1\cdots A_t]\leq b\). Then
\(|G|\leq b\prod_j a_j\). If a finite quotient of \(G\) has this latter
order, its quotient map is an isomorphism.

\emph{Proof.} Normality makes each partial product a subgroup. The
multiplication map from \(A_1\times\cdots\times A_t\) onto the product
is surjective as a map of sets, so that product has order at most
\(\prod_j a_j\). Multiplication by the quotient index proves the bound
and finiteness of \(G\). A quotient attaining the bound forces equality
of the finite orders and hence a trivial kernel. \(\square\)

\textbf{Corollary UA.Search (rank and excess classify the finite
search).} Let T be a specified forest on d vertices and m edges, and let
k≥0. Set

\[
 F_T=\langle x_i\mid x_i^2=1,\ [x_i,x_j]=1\text{ for nonedges}\rangle,
 \quad U_{T,k}=F_T/\gamma_{k+3}(F_T),\quad K=\gamma_3(U_{T,k}).
\]

The ordered IP markings by nonidentity involutions with graph T and
order \(2^{d+m+k}\) are exactly the quotients \(U_{T,k}/N\), where N is
normal, contained in K, \([K:N]=2^k\), and the images
\((N(K\cap(U_{T,k})_J))/N\) have IP for every pair J,L. Generator-set
equivalence is exactly the action of Aut(T) on these kernels.

\emph{Proof.} Killing any set of generators preserves every defining
relator of \(F_T\). Collection in its class-two quotient gives at most
d+m binary coordinates, and (F0) realizes all of them, so this quotient
is precisely the marked model \(Q_T\). Every target has that same
quotient by UA.Forest. These are exactly UA.Nilpotent's hypotheses. Its
all-parabolic condition and UA.Ladder give the stated bijection and
exhaustive search. An ordered isomorphism must commute with the
generating map, hence has the same kernel. A permutation of the
generator set preserves the noncommutation graph and extends to an
automorphism of \(F_T\) and \(U_{T,k}\); this proves the last claim. ∎

The last corollary is a finite classification procedure, with the
marking retained. First form the finite group \(U_{T,k}\); then iterate
UA.Ladder inside \(\gamma_3(U_{T,k})\), test the explicit intersection
condition of UA.Nilpotent, and take the declared graph-automorphism
orbits. The theorem proves that this procedure terminates and misses no
target marking. Any reported numbers of groups, named presentations or
universal orders still require execution and an application-specific
identification. No numerical census is assumed or reported in this
paper.
